\section{Higher dimensions and open problems}
\label{sec:higher}
The clique formula of Section~\ref{sec:clique} holds in all dimensions,
while the exact chromatic results of Sections~\ref{sec:line} and
\ref{sec:chromatic} concern dimension six. This section records a
reduction for higher-dimensional coloring problems. We first give a degree-versus-list deletion
rule that reduces a \(k\)-coloring problem in dimension \(n\) to a
smaller residual instance, without changing \(k\)-colorability.
We then apply it to \(n=7\), where a \(16\)-clique is available and
the rule removes all vertices of dimension at least four, leaving a
residual instance of \(14\,589\) vertices and \(953\,057\) edges.
Finally, we state the open problems that remain: the exact
chromatic number for \(n\ge7\), the question of whether
\(\chi>\omega\) holds in infinitely many dimensions, and the structural
question of whether an optimal \(15\)-coloring in dimension six
can be described without finite computation.
\begin{proposition}\label{prop:dimension-deletion}
Let \(K\) be a fixed \(k\)-clique in \(\OG_n^*\).  If \(N(n-d)<k\),
every uncolored \(d\)-dimensional vertex can be deleted by the
degree-versus-list rule.  Consequently, \(k\)-colorability of the
residual instance after such deletions is equivalent to
\(k\)-colorability of the full graph.
\end{proposition}
\begin{proof}
For a \(d\)-space \(U\), the orthogonal complement has dimension
\(n-d\).  Every neighbor of \(U\) is a nonzero subspace of
\(U^\perp\), so \(\deg U\leq N(n-d)\).  If \(f\) fixed clique
vertices are adjacent to \(U\), then its residual degree is at most
\(N(n-d)-f\), while its list size is \(k-f\).  The stated inequality
therefore guarantees an available color on restoring \(U\).  Apply
this argument in reverse deletion order.
\end{proof}

For \(n=7\), the nonzero subspaces of \(T=\Span(3,12,48)\) together
with \(\Span(64)\) give a \(16\)-clique.  Since \(N(3)=15<16\),
Proposition~\ref{prop:dimension-deletion} removes all \(14\,606\)
vertices of dimension at least four from the \(29\,211\)-vertex
graph.  The remaining uncolored instance has \(14\,589\) vertices
and \(953\,057\) edges.  The direct pairwise list-color encoding
would have \(214\,355\) variables and \(12\,984\,679\) clauses.
These exact instance sizes guide further geometric reduction and
resource-bounded search.  They are measured preprocessing counts;
no \(16\)-coloring or checked nonexistence certificate for
dimension seven is supplied.

Problem~4.3 remains open for \(n\geq7\), as does Problem~4.4 of
\cite{Nikandish2026}.  Within dimension six, a structural description
of an optimal full-subspace coloring remains an immediate question.

The negative results of Section~\ref{sec:geometry} --- in particular
Proposition~\ref{prop:lift-features} and
Proposition~\ref{prop:stabT-orbit} --- rule out colorings depending
only on the exact profiles or the specified lift features, as well as
colorings invariant under the stabilizer of \(T\).
A structural proof of
\(\chi(\OG_6^*)=15\) must therefore use auxiliary data that break
this symmetry and retain enough placement information to separate
adjacent vertices.

\begin{problem}\label{prob:structural}
Can one specify auxiliary geometric data and a nontrivial proper subgroup
\(H<\operatorname{Stab}_{O_6(\F)}(T)\), and construct an explicit proper
\(15\)-coloring \(c\) from those data such that
\(c(hU)=h\cdot c(U)\) for \(h\in H\), where the color labels are the
nonzero subspaces of \(T\) with their induced \(H\)-action?

The construction should prove availability and separation without
looking up the archived coloring certificate. Such a construction would
give a conceptual proof of \(\chi(\OG_6^*)=15\).
The checked \(12\)-coloring of \(\Gamma_6\) establishes its chromatic number;
equivariance under the full orthogonal group is not established here.
\end{problem}

\begin{remark}\label{prob:gamma-six}
The line-graph question is settled by Theorem~\ref{thm:lines}:
\(\chi(\Gamma_6)=12\), whereas \(\chi(\OG_6^*)=15\).
Thus passing from lines to all nonzero subspaces increases the
chromatic number by exactly three in dimension six.
\end{remark}
